\documentclass{cumcmthesis}
%\documentclass[withoutpreface,bwprint]{cumcmthesis}

\title{回焊炉炉温曲线的机理建模、参数辨识与多目标优化}
\tihao{A}
\baominghao{4321}
\schoolname{上海电机学院}
\membera{黄张}
\memberb{李雨树}
\memberc{张天毅}
\supervisor{武文佳}
\yearinput{2026}
\monthinput{08}
\dayinput{15}

\begin{document}
\maketitle

\begin{abstract}
针对2020年高教社杯全国大学生数学建模竞赛A题“炉温曲线”，本文从回焊炉内焊接区域的受热机理出发，建立“炉内环境温度场—工件温度响应—工艺指标—参数优化”的统一数学模型。

首先，将炉膛沿传送方向划分为加热温区、温区间隙及炉前炉后过渡区域，并根据题目给出的温区长度、间隙长度和设定温度构造分段环境温度函数。考虑焊接区域厚度较小，采用集总热容思想，以牛顿冷却定律描述工件与炉内环境之间的换热过程，建立一阶非齐次常微分方程。利用附件实验数据，通过最小二乘准则辨识综合换热参数，从而完成模型参数标定，并在问题一给定工况下预测焊接区域中心温度曲线及指定位置温度。

其次，针对问题二，将传送带速度作为唯一决策变量，以峰值温度、升温斜率、降温斜率、150--190℃升温时间以及高于217℃持续时间为约束，建立最大过炉速度的一维约束优化模型。针对速度增大导致工件受热时间缩短的机理，先确定可行速度区间，再采用区间搜索与二分细化确定最大可行速度，并通过边界复核识别活跃约束。

再次，针对问题三，以“217℃以上至峰值之间的超温面积最小”为优化目标，以四类温区设定温度和传送带速度为决策变量，在全部制程界限下构建非线性约束优化模型。采用“差分进化全局搜索+SLSQP局部精化”的混合策略提高求解稳定性，并通过多组初值和约束复核降低局部最优风险。

最后，针对问题四，在问题三目标基础上引入峰值两侧温度曲线的对称性指标，构造归一化时间坐标下的曲线差异度量，并采用归一化加权方法形成双目标优化模型。通过比较问题三与问题四的面积、对称性及制程指标，分析面积最优与曲线对称之间的权衡关系。

本文的主要贡献在于：建立了可解释的热过程机理模型；将四个问题统一到同一参数化模型框架中；对优化问题给出明确的可行域、活跃约束和数值求解流程；同时通过敏感性分析与边界复核讨论模型稳定性。由于当前稿件未附原始实验数据及最终程序运行结果，所有必须由数据计算得到的数值结果均明确标记为待由支撑程序填入，避免人为虚构竞赛结果。
\keywords{炉温曲线\quad 集总热容\quad 牛顿冷却定律\quad 参数辨识\quad 约束优化\quad 差分进化}
\end{abstract}

\tableofcontents

\section{问题重述}

回焊炉是表面贴装技术生产线中的关键设备。电路板随传送带以近似恒定速度通过炉膛，在不同温区中经历预热、升温、回流和冷却等过程。焊接区域的中心温度随时间变化形成炉温曲线，其峰值温度、升降温速率以及在特定温度区间内的停留时间直接影响焊点质量。

题目给出了11个小温区的长度、相邻温区之间的间隙、炉前炉后长度以及各温区的设定温度，并给出了表~\ref{tab:process} 所示的五项制程界限。要求建立焊接区域温度变化模型，并进一步完成以下四个问题。

\begin{enumerate}
    \item \textbf{问题一：} 在给定温区设定温度和传送带速度下，建立焊接区域中心温度模型，预测完整炉温曲线，并计算温区3中点、温区6中点、温区7中点以及温区8结束处的温度。
    \item \textbf{问题二：} 固定温区设定温度，在传送带速度可调范围内，求满足全部制程界限的最大允许过炉速度。
    \item \textbf{问题三：} 在规定温度调节范围和速度范围内，以217℃以上至峰值之间的超温面积最小为目标，寻找最优温区温度组合和传送带速度。
    \item \textbf{问题四：} 在问题三的基础上进一步要求峰值两侧曲线尽量对称，研究超温面积和曲线对称性之间的权衡，并求得综合最优方案。
\end{enumerate}

\begin{table}[H]
\centering
\caption{制程界限}
\label{tab:process}
\begin{tabular}{cccc}
\toprule
界限名称 & 最低值 & 最高值 & 单位\\
\midrule
温度上升斜率 & 0 & 3 & ${}^\circ$C/s\\
温度下降斜率 & -3 & 0 & ${}^\circ$C/s\\
150--190℃升温时间 & 60 & 120 & s\\
大于217℃持续时间 & 40 & 90 & s\\
峰值温度 & 240 & 250 & ${}^\circ$C\\
\bottomrule
\end{tabular}
\end{table}

\section{问题分析}

\subsection{总体建模思路}

四个问题具有明显的递进关系。问题一解决温度响应模型及参数辨识问题；问题二在模型固定后仅改变传送带速度；问题三进一步将温区设定温度纳入决策变量；问题四在问题三基础上增加曲线形态指标。

因此，本文采用统一的计算框架：
\[
\boxed{
\text{环境温度场}
\longrightarrow
\text{温度ODE}
\longrightarrow
\text{制程指标}
\longrightarrow
\text{约束判定}
\longrightarrow
\text{优化决策}
}
\]

其中，环境温度场由空间位置决定，工件温度由ODE响应，五项制程指标由温度曲线计算得到，优化问题则以这些指标作为约束或目标。

\subsection{问题一分析}

问题一的核心是确定焊接区域温度与炉内环境之间的动态关系。由于工件很薄，可以首先检验集总热容假设；若Biot数足够小，则无需建立工件内部空间温度场，而只需跟踪其整体温度。

\subsection{问题二分析}

问题二只有速度一个决策变量。速度增大意味着工件在每个温区中的停留时间缩短，因此通常会导致峰值温度、150--190℃升温时间以及217℃以上持续时间下降。于是问题二可以转化为一维可行域右端点搜索问题。

需要强调的是：这种单调关系属于模型机理和数值实验共同支持的性质，而不能仅凭直觉直接作为数学定理使用。因此，实际求解中应对指标随速度变化的曲线进行数值验证，再使用二分法。

\subsection{问题三分析}

问题三同时调节温区温度和传送带速度，决策空间维数明显增加，目标函数又通过ODE间接依赖于决策变量，因此属于非线性、非凸、带路径型制程约束的优化问题。采用全局启发式算法寻找候选解，再利用局部约束优化算法精化，可以兼顾搜索范围和局部精度。

\subsection{问题四分析}

问题四中的“对称”不是简单比较峰值前后两个时刻的温度，而需要消除峰值前后时间跨度不同带来的尺度影响。为此，在超过217℃的峰前、峰后区间内分别进行归一化时间映射，再计算两段温度曲线之间的差异。

\section{模型假设}

\begin{enumerate}
    \item 炉内温度场达到准稳态，即在研究时间尺度内，环境温度主要随炉长方向位置变化，可表示为 $T_a(x)$。
    \item 传送带速度在一次过炉过程中保持恒定，因此焊接区域位置满足 $x(t)=vt$。
    \item 焊接区域内部温度近似均匀，可采用集总热容模型描述。
    \item 工件与炉内环境之间的换热可由牛顿冷却定律描述，辐射等效应并入综合换热参数。
    \item 由附件实验数据辨识得到的综合换热参数在其他题设工况下保持不变。
    \item 温区间隙内的环境温度采用相邻温区设定值的线性过渡进行近似。
    \item 优化过程中边界约束允许取等号，即满足制程界限边界即可视为可行。
\end{enumerate}

其中，第5、6条是模型外推的关键假设。其合理性将在模型评价部分通过敏感性分析和边界复核进一步讨论。

\section{符号说明}

\begin{center}
\begin{tabular}{ccc}
\toprule
符号 & 含义 & 单位\\
\midrule
$x$ & 炉长方向位置 & cm\\
$t$ & 时间 & s\\
$v$ & 传送带速度 & cm/s\\
$T_a(x)$ & 炉内环境温度 & ${}^\circ$C\\
$T(t)$ & 焊接区域中心温度 & ${}^\circ$C\\
$\alpha$ & 综合换热参数 & s$^{-1}$\\
$L$ & 单个小温区长度 & cm\\
$g$ & 相邻温区间隙长度 & cm\\
$P$ & 炉前或炉后长度 & cm\\
$T_{set,i}$ & 第$i$温区设定温度 & ${}^\circ$C\\
$T_{pk}$ & 峰值温度 & ${}^\circ$C\\
$t_{217}^{in}$ & 升温时首次达到217℃的时刻 & s\\
$t_{217}^{out}$ & 降温时最后高于217℃的时刻 & s\\
$\Delta t_{>217}$ & 大于217℃持续时间 & s\\
$\Delta t_{150\sim190}$ & 150--190℃升温时间 & s\\
$S_r$ & 最大升温斜率 & ${}^\circ$C/s\\
$S_d$ & 最小降温斜率 & ${}^\circ$C/s\\
$A$ & 217℃至峰值的超温面积 & ${}^\circ$C$\cdot$s\\
$D$ & 曲线对称性差异指标 & ${}^\circ$C\\
$\boldsymbol u$ & 优化决策向量 & --\\
\bottomrule
\end{tabular}
\end{center}

\section{问题一：炉温曲线模型与参数辨识}

\subsection{炉内环境温度场}

题目给出单个温区长度为30.5 cm，相邻温区之间间隙为5 cm，炉前和炉后均为25 cm。因此第$i$个温区起点和终点分别为
\begin{equation}
x_i^s=25+35.5(i-1),\qquad
x_i^e=x_i^s+30.5,\qquad i=1,\ldots,11.
\end{equation}

炉膛总长度为
\begin{equation}
X_{\rm total}
=25+11\times30.5+10\times5+25
=435.5\ {\rm cm}.
\end{equation}

设第$i$温区设定温度为 $T_{set,i}$。在温区内部取环境温度为设定值，在相邻温区之间采用线性过渡，在炉前和炉后区域同样采用线性过渡，则
\begin{equation}
T_a(x)=
\begin{cases}
25+\dfrac{x}{25}(T_{set,1}-25),
&0\le x<25,\\[8pt]
T_{set,i},
&x_i^s\le x\le x_i^e,\\[8pt]
T_{set,i}
+\dfrac{x-x_i^e}{x_{i+1}^s-x_i^e}
(T_{set,i+1}-T_{set,i}),
&x_i^e<x<x_{i+1}^s,\\[8pt]
T_{set,11}
+\dfrac{x-x_{11}^e}{25}(25-T_{set,11}),
&x_{11}^e\le x\le435.5.
\end{cases}
\label{eq:Ta}
\end{equation}

该处理保证了环境温度函数在各过渡位置连续，避免数值积分时出现无物理意义的瞬时无限温差。

\subsection{集总热容模型}

焊接区域厚度较小，采用集总热容模型。由能量守恒，
\begin{equation}
\rho c_p V\frac{{\rm d}T}{{\rm d}t}
=hS\left[T_a(x(t))-T(t)\right],
\end{equation}
其中 $S$ 为有效换热面积。令
\begin{equation}
\alpha=\frac{hS}{\rho c_pV},
\end{equation}
并利用 $x(t)=vt$，得到
\begin{equation}
\boxed{
\frac{{\rm d}T}{{\rm d}t}
=\alpha\left[T_a(vt)-T(t)\right]
},
\qquad
T(0)=25.
\label{eq:ode}
\end{equation}

该模型为一阶非齐次线性常微分方程。由于 $T_a(vt)$ 为分段函数，本文采用四阶Runge--Kutta方法进行数值积分。

\subsection{集总热容假设的合理性}

集总热容法的适用性可由Biot数判断：
\begin{equation}
Bi=\frac{hL_c}{k}.
\end{equation}
根据题设工件厚度0.15 mm，若取特征长度
\[
L_c=\frac{0.15}{2}\ {\rm mm}=7.5\times10^{-5}\ {\rm m},
\]
并采用与材料及换热环境相符的参数量级进行估算，可得到 $Bi<0.1$。因此工件内部导热时间尺度明显小于其与环境换热的时间尺度，采用单温度集总模型具有合理性。

\subsection{综合换热参数辨识}

附件给出了实验工况下的观测温度数据。对于候选参数 $\alpha$，由式~\eqref{eq:ode}获得模型温度 $T_{\rm model}(t_k;\alpha)$，以观测值 $T_{\rm obs}(t_k)$ 为基准建立最小二乘目标：
\begin{equation}
\alpha^*
=
\arg\min_{\alpha>0}
\sum_{k=1}^{n}
\left[
T_{\rm model}(t_k;\alpha)-T_{\rm obs}(t_k)
\right]^2.
\label{eq:alpha}
\end{equation}

为避免只报告一个拟合指标，本文同时采用决定系数 $R^2$ 和均方根误差RMSE评价模型：
\begin{equation}
R^2
=
1-
\frac{\sum_k(T_{\rm obs,k}-T_{\rm model,k})^2}
{\sum_k(T_{\rm obs,k}-\bar T_{\rm obs})^2},
\end{equation}
\begin{equation}
RMSE
=
\sqrt{
\frac{1}{n}
\sum_k
(T_{\rm obs,k}-T_{\rm model,k})^2
}.
\end{equation}

\textbf{说明：当前上传材料未包含附件原始实验数据，因此本节的 $\alpha^*$、$R^2$ 和RMSE不能凭空填写。最终竞赛稿必须由支撑程序根据原始数据自动生成。}

\subsection{问题一给定工况}

问题一的设定温度为
\[
(T_{1\sim5},T_6,T_7,T_{8\sim9},T_{10\sim11})
=(173,198,230,257,25)\ ^\circ{\rm C},
\]
传送带速度为
\[
v=78\ {\rm cm/min}
=1.3\ {\rm cm/s}.
\]

四个指定位置为
\[
x_{3,\rm mid}=111.25\ {\rm cm},
\quad
x_{6,\rm mid}=217.75\ {\rm cm},
\]
\[
x_{7,\rm mid}=253.25\ {\rm cm},
\quad
x_{8,\rm end}=304.00\ {\rm cm}.
\]

对应时刻为
\begin{equation}
t=\frac{x}{v},
\end{equation}
即约为85.58 s、167.50 s、194.81 s和233.85 s。

\begin{table}[H]
\centering
\caption{问题一指定位置计算结果}
\label{tab:q1result}
\begin{tabular}{cccc}
\toprule
位置 & 坐标/cm & 时刻/s & 温度/${}^\circ$C\\
\midrule
温区3中点 & 111.25 & 85.58 & 待程序填入\\
温区6中点 & 217.75 & 167.50 & 待程序填入\\
温区7中点 & 253.25 & 194.81 & 待程序填入\\
温区8结束处 & 304.00 & 233.85 & 待程序填入\\
\bottomrule
\end{tabular}
\end{table}

\subsection{数值误差控制}

采用RK4计算时，内部积分步长记为 $\Delta t$，最终按0.5 s输出。为验证步长对结果的影响，建议分别采用
\[
\Delta t=1.0,\quad0.5,\quad0.25,\quad0.125\ {\rm s}
\]
进行计算，并比较峰值温度和各关键制程指标。当连续两次细化后主要结果变化小于预设容差时，将对应步长作为正式计算步长。

这一处理比直接给出“0.5 s采样误差约为某个固定温度值”更加严谨，因为ODE离散误差与输出采样误差是两个不同来源。

\section{问题二：最大允许过炉速度}

\subsection{制程指标定义}

给定传送带速度 $v$，由温度曲线计算五项指标：
\begin{align}
T_{pk}(v)&=\max_t T(t),\\
S_r(v)&=\max_{t:T'(t)>0}T'(t),\\
S_d(v)&=\min_{t:T'(t)<0}T'(t),\\
\Delta t_{150\sim190}(v)
&=
t_{190}^{\uparrow}-t_{150}^{\uparrow},\\
\Delta t_{>217}(v)
&=
t_{217}^{\rm out}-t_{217}^{\rm in}.
\end{align}

其中上标$\uparrow$表示升温阶段，$t_{217}^{\rm in}$和$t_{217}^{\rm out}$分别表示曲线首次进入和最后离开217℃以上区域的时刻。

\subsection{约束优化模型}

题目给定问题二温区设定温度
\[
(T_{1\sim5},T_6,T_7,T_{8\sim9},T_{10\sim11})
=(182,203,237,254,25)\ ^\circ{\rm C},
\]
速度范围为
\[
65\le v\le100\ {\rm cm/min}.
\]

建立优化模型
\begin{equation}
\begin{aligned}
\max\quad &v\\
{\rm s.t.}\quad
&65\le v\le100,\\
&240\le T_{pk}(v)\le250,\\
&0\le S_r(v)\le3,\\
&-3\le S_d(v)\le0,\\
&60\le\Delta t_{150\sim190}(v)\le120,\\
&40\le\Delta t_{>217}(v)\le90.
\end{aligned}
\label{eq:q2opt}
\end{equation}

定义可行集
\begin{equation}
\mathcal V=
\left\{
v\in[65,100]:
\text{式~\eqref{eq:q2opt}全部约束成立}
\right\},
\end{equation}
则
\begin{equation}
v_{\max}=\max\mathcal V.
\end{equation}

\subsection{搜索算法}

由于只有一个决策变量，采用“粗扫描+二分细化”的方法。

首先以1 cm/min为间隔扫描 $[65,100]$，得到可行速度区间的大致位置。随后，对可行区间右端点附近进行二分搜索。对于每个候选速度，均完整求解温度ODE并重新计算五项制程指标，而不是仅用某一个指标近似判断。

二分停止条件取
\[
|v_R-v_L|<10^{-3}\ {\rm cm/min}.
\]

若数值实验验证 $T_{pk}$、$\Delta t_{150\sim190}$和$\Delta t_{>217}$在研究区间内随速度单调下降，则最大速度可由相应下限约束的临界速度确定：
\begin{equation}
v_{\max}
=
\min
\left\{
v_{T_{pk}=240},
v_{\Delta t_{150\sim190}=60},
v_{\Delta t_{>217}=40}
\right\},
\label{eq:q2active}
\end{equation}
但式~\eqref{eq:q2active}只有在单调性得到数值验证后才使用。

\subsection{结果与活跃约束}

\begin{table}[H]
\centering
\caption{问题二最大速度及制程指标}
\label{tab:q2result}
\begin{tabular}{ccc}
\toprule
指标 & $v_{\max}$处结果 & 允许范围\\
\midrule
最大速度/cm/min & 待程序填入 & [65,100]\\
峰值温度/${}^\circ$C & 待程序填入 & [240,250]\\
最大升温斜率/${}^\circ$C/s & 待程序填入 & [0,3]\\
最小降温斜率/${}^\circ$C/s & 待程序填入 & [-3,0]\\
150--190℃时间/s & 待程序填入 & [60,120]\\
217℃以上时间/s & 待程序填入 & [40,90]\\
\bottomrule
\end{tabular}
\end{table}

最终应将“关键约束”定义为在 $v_{\max}$ 处具有最小归一化裕量的约束，而不应在没有计算结果的情况下预先断言某一个约束必然活跃。

进一步取略高于最优速度的 $v_{\max}+\delta$，检查其至少有一项制程指标越界，从而验证 $v_{\max}$ 的确为可行域右端点。

\section{问题三：超温面积最小化}

\subsection{决策变量}

问题三允许温区设定温度在规定范围内调整。将1--5区、6区、7区、8--9区分别作为四组温度变量，则
\begin{equation}
\boldsymbol u=
(T_{1\sim5},T_6,T_7,T_{8\sim9},v).
\end{equation}

其取值范围为
\begin{equation}
\begin{aligned}
165&\le T_{1\sim5}\le185,\\
185&\le T_6\le205,\\
225&\le T_7\le245,\\
245&\le T_{8\sim9}\le265,\\
65&\le v\le100.
\end{aligned}
\label{eq:q3bounds}
\end{equation}

\subsection{目标函数的严格定义}

题目要求最小化“超过217℃至峰值”的面积。因此不能将整个过炉时间内的 $(T-217)$ 直接积分。设
\[
t_{217}^{\uparrow}<t_{pk}<t_{217}^{\downarrow},
\]
其中 $t_{217}^{\uparrow}$ 为升温过程中首次达到217℃的时刻，$t_{pk}$ 为峰值时刻，则本文将问题三目标定义为
\begin{equation}
\boxed{
A(\boldsymbol u)
=
\int_{t_{217}^{\uparrow}}^{t_{pk}}
[T(t;\boldsymbol u)-217]\,{\rm d}t
}.
\label{eq:area}
\end{equation}

该定义严格对应“217℃以上至峰值”的升温侧面积，避免把降温段面积误计入目标函数。

\subsection{约束条件}

优化过程中必须同时满足全部制程界限：
\begin{equation}
\begin{cases}
240\le T_{pk}(\boldsymbol u)\le250,\\
0\le S_r(\boldsymbol u)\le3,\\
-3\le S_d(\boldsymbol u)\le0,\\
60\le\Delta t_{150\sim190}(\boldsymbol u)\le120,\\
40\le\Delta t_{>217}(\boldsymbol u)\le90.
\end{cases}
\label{eq:q3constraints}
\end{equation}

因此问题三为
\begin{equation}
\boxed{
\min_{\boldsymbol u}A(\boldsymbol u)
}
\end{equation}
subject to 式~\eqref{eq:q3bounds}和~\eqref{eq:q3constraints}。

\subsection{数值求解}

由于目标函数通过ODE间接获得，且约束边界可能造成可行域非凸，本文采用两阶段混合优化：

\begin{enumerate}
    \item 使用差分进化算法在整个决策空间内进行全局搜索，获得若干候选解；
    \item 对候选解采用SLSQP进行局部约束精化；
    \item 对最终解重新以高精度ODE求解器计算所有制程指标；
    \item 对不同随机种子重复运行，比较最终目标函数值和可行性；
    \item 只有同时满足全部约束的解才进入最终候选集。
\end{enumerate}

与单纯使用罚函数相比，上述流程将“搜索”与“最终可行性判定”分开，更适合竞赛论文中对结果可靠性的论证。

\subsection{问题三结果}

\begin{table}[H]
\centering
\caption{问题三最优参数及制程指标}
\label{tab:q3result}
\begin{tabular}{ccc}
\toprule
变量/指标 & 最优值 & 允许范围\\
\midrule
$T_{1\sim5}$/℃ & 待程序填入 & [165,185]\\
$T_6$/℃ & 待程序填入 & [185,205]\\
$T_7$/℃ & 待程序填入 & [225,245]\\
$T_{8\sim9}$/℃ & 待程序填入 & [245,265]\\
$v$/cm/min & 待程序填入 & [65,100]\\
\midrule
$A^*/({}^\circ{\rm C}\cdot{\rm s})$ & 待程序填入 & 越小越优\\
$T_{pk}$/℃ & 待程序填入 & [240,250]\\
$\Delta t_{>217}$/s & 待程序填入 & [40,90]\\
$\Delta t_{150\sim190}$/s & 待程序填入 & [60,120]\\
$S_r$/℃/s & 待程序填入 & [0,3]\\
$S_d$/℃/s & 待程序填入 & [-3,0]\\
\bottomrule
\end{tabular}
\end{table}

在最终论文中，应根据表~\ref{tab:q3result} 判断哪些约束在最优解处处于边界，并给出相应的物理解释。例如，若峰值温度接近240℃，说明为了减小超温面积，优化过程倾向于将回流热负荷压缩到满足最低工艺要求的程度。

\section{问题四：兼顾超温面积与曲线对称性}

\subsection{对称性指标}

为了评价峰值两侧的对称程度，先定义217℃以上区域的四个关键时刻：
\[
t_{217}^{\uparrow}<t_{pk}<t_{217}^{\downarrow}.
\]

将峰值前的区间
\[
[t_{217}^{\uparrow},t_{pk}]
\]
和峰值后的区间
\[
[t_{pk},t_{217}^{\downarrow}]
\]
分别映射到归一化坐标 $\theta\in[0,1]$：
\begin{equation}
T_{\rm pre}(\theta)
=
T\left(
t_{pk}-\theta(t_{pk}-t_{217}^{\uparrow})
\right),
\end{equation}
\begin{equation}
T_{\rm post}(\theta)
=
T\left(
t_{pk}+\theta(t_{217}^{\downarrow}-t_{pk})
\right).
\end{equation}

定义对称性差异指标
\begin{equation}
\boxed{
D(\boldsymbol u)
=
\int_0^1
\left|
T_{\rm pre}(\theta)-T_{\rm post}(\theta)
\right|
\,{\rm d}\theta
}.
\label{eq:sym}
\end{equation}

$D$越小表示归一化后的峰前、峰后曲线越接近。

该指标的优点是消除了峰值两侧持续时间不同造成的直接尺度影响；其局限是它反映的是“形状对称”，而不是原始时间尺度上的绝对对称。

\subsection{双目标处理}

问题四同时要求减小超温面积 $A$ 和对称性差异 $D$。二者量纲不同，因此不能直接线性相加。本文采用归一化加权方法：
\begin{equation}
F(\boldsymbol u)
=
(1-\lambda)
\frac{A(\boldsymbol u)}{A_{\rm ref}}
+
\lambda
\frac{D(\boldsymbol u)}{D_{\rm ref}},
\qquad
0\le\lambda\le1.
\label{eq:q4obj}
\end{equation}

其中 $A_{\rm ref}$ 和 $D_{\rm ref}$ 为参考尺度。为避免参考尺度依赖单个偶然解，正式计算中建议采用问题三最优解及多次可行初始解的统计尺度进行归一化；若采用问题三最优解作为参考，则必须在论文中明确说明。

本文主方案取
\[
\lambda=0.4,
\]
并通过改变 $\lambda$ 进行敏感性分析。若$\lambda$变化较小而最优参数和主要制程指标保持稳定，则说明结论具有一定鲁棒性。

\subsection{问题四求解}

沿用问题三的DE+SLSQP混合求解框架，只将目标函数替换为式~\eqref{eq:q4obj}。所有制程界限保持不变。

\begin{table}[H]
\centering
\caption{问题四与问题三优化结果对比}
\label{tab:q4result}
\begin{tabular}{cccc}
\toprule
变量/指标 & 问题四 & 问题三 & 相对变化\\
\midrule
$T_{1\sim5}$/℃ & 待程序填入 & 待程序填入 & 待程序填入\\
$T_6$/℃ & 待程序填入 & 待程序填入 & 待程序填入\\
$T_7$/℃ & 待程序填入 & 待程序填入 & 待程序填入\\
$T_{8\sim9}$/℃ & 待程序填入 & 待程序填入 & 待程序填入\\
$v$/cm/min & 待程序填入 & 待程序填入 & 待程序填入\\
\midrule
$A^*/({}^\circ{\rm C}\cdot{\rm s})$ & 待程序填入 & 待程序填入 & 待程序填入\\
$D^*/{}^\circ{\rm C}$ & 待程序填入 & 待程序填入 & 待程序填入\\
$T_{pk}$/℃ & 待程序填入 & 待程序填入 & 待程序填入\\
$\Delta t_{>217}$/s & 待程序填入 & 待程序填入 & 待程序填入\\
\bottomrule
\end{tabular}
\end{table}

最终应报告
\begin{equation}
\frac{A_4-A_3}{A_3}\times100\%
\end{equation}
和
\begin{equation}
\frac{D_4-D_3}{D_3}\times100\%
\end{equation}
以量化面积与对称性之间的权衡。

\section{模型验证与敏感性分析}

\subsection{参数敏感性}

综合换热参数$\alpha$是整个模型最重要的标定参数。令
\[
\alpha_\eta=(1+\eta)\alpha^*,
\]
其中 $\eta$ 可取 $-10\%,-5\%,0,5\%,10\%$。分别重新计算问题二至问题四的结果，观察$v_{\max}$、$T_{pk}$和目标函数的变化。

若参数扰动下最优结论变化较小，则说明模型对参数误差具有一定鲁棒性；若变化明显，则应在论文中明确指出$\alpha$辨识精度是模型主要误差来源。

\subsection{数值步长敏感性}

分别采用不同RK4步长计算问题一至问题四的主要指标，并定义相对变化率
\begin{equation}
E_q
=
\frac{|q_{\Delta t}-q_{\Delta t/2}|}
{\max(1,|q_{\Delta t/2}|)}.
\end{equation}

当 $E_q$ 足够小时，认为数值离散误差不会影响最优方案的主要结论。

\subsection{边界复核}

对于每一个最终优化结果，重新使用独立的高精度计算流程验证：
\[
240\le T_{pk}\le250,
\quad
0\le S_r\le3,
\quad
-3\le S_d\le0,
\]
\[
60\le\Delta t_{150\sim190}\le120,
\quad
40\le\Delta t_{>217}\le90.
\]

对于问题二，还应验证
\[
v_{\max}-\varepsilon
\]
可行，而
\[
v_{\max}+\varepsilon
\]
不可行，其中 $\varepsilon$ 取一个足够小的速度扰动。这比仅报告一个速度数值更能证明“最大允许速度”的结论。

% =========================================================
% 十、模型评价
% =========================================================

\section{模型评价}

\subsection{模型优点}

\begin{enumerate}
    \item \textbf{物理意义明确。}
    
    本文以焊接区域与炉内环境之间的动态换热为基础，将复杂的温度变化过程抽象为一阶常微分方程。模型中的环境温度、传送带速度、工件温度以及综合换热参数均具有明确的物理含义，能够较好地反映工件在不同温区中的受热与冷却过程。

    \item \textbf{四个问题具有统一的模型基础。}
    
    四个问题均由同一炉温响应模型产生温度曲线，再由温度曲线计算峰值温度、升降温斜率以及特定温度区间内的停留时间等制程指标。在此基础上，问题二进一步调整传送带速度，问题三同时调整温区温度和传送带速度，问题四进一步考虑峰值两侧曲线的对称性，从而形成由预测到优化逐层递进的统一建模框架。

    \item \textbf{优化变量少而解释性强。}
    
    对复杂炉温系统进行建模时，将大量传热因素统一归入综合换热参数 $\alpha$，避免引入难以辨识的过多参数。在已有实验数据条件下，该处理能够在模型复杂度和计算稳定性之间取得平衡。

    \item \textbf{数值求解方法与问题结构相匹配。}
    
    问题二仅包含一个决策变量，因此采用粗搜索与二分细化能够有效确定最大可行速度；问题三和问题四属于多变量非线性约束优化问题，因此采用全局搜索与局部精化相结合的方法，可以提高搜索稳定性与最终解的精度。

    \item \textbf{结果具有可验证性。}
    
    最终优化结果不仅需要满足全部制程约束，还可通过改变积分步长、扰动综合换热参数以及重复运行优化算法等方式进行稳定性检验，从而降低数值误差和随机优化误差对结论的影响。
\end{enumerate}


\subsection{模型不足}

\begin{enumerate}
    \item \textbf{综合换热参数采用常数近似。}
    
    实际回焊过程中，对流与辐射换热强度可能随温度、位置以及工件状态发生变化。本文将其统一等效为参数 $\alpha$，虽然能够降低模型复杂度，但会带来一定的模型偏差。

    \item \textbf{温区间隙温度采用工程近似。}
    
    由于缺乏温区间隙处的实际炉膛温度测量数据，本文采用相邻温区设定温度之间的线性过渡进行近似。如果能够获得实际炉内温场数据，则可以进一步利用插值或实验标定提高环境温度函数的准确性。

\end{enumerate}


\subsection{模型改进方向}

若进一步提高模型精度，可以从以下三个方向进行改进。

首先，可将综合换热参数由常数 $\alpha$ 扩展为温度或位置的函数，例如
\[
\alpha=\alpha(T,x),
\]
从而更加准确地刻画不同炉段的实际换热特性。

其次，可以利用实际炉膛温度测量数据替代温区间隙处的线性插值，使环境温度函数更接近真实炉温场。

最后，可以在优化过程中进一步引入设备能耗、温区调节幅度以及工艺鲁棒性等因素，将单纯的参数优化扩展为更加符合工业实际的多目标优化。


% =========================================================
% 十一、结论
% =========================================================

\section{结论}

本文针对2020年高教社杯全国大学生数学建模竞赛A题“炉温曲线”，建立了由环境温度场、温度响应模型、制程指标计算以及参数优化组成的统一数学模型。

\subsection{问题一结论}

首先，根据炉膛中各温区的位置、长度和设定温度构造了炉内环境温度函数 $T_a(x)$，并利用传送带匀速运动关系
\[
x(t)=vt
\]
将空间温度场转化为时间函数。

在集总热容假设和牛顿冷却定律基础上，建立
\[
\frac{{\rm d}T}{{\rm d}t}
=
\alpha
\left[
T_a(vt)-T(t)
\right]
\]
的一阶常微分方程，并采用数值积分方法获得焊接区域完整温度曲线。

进一步根据实验观测数据，以最小二乘准则辨识综合换热参数 $\alpha$，从而建立可用于后续问题计算的温度预测模型。

\subsection{问题二结论}

问题二以传送带速度 $v$ 为唯一决策变量，在全部制程界限下建立最大允许过炉速度模型：
\[
\max v.
\]

通过对速度区间进行粗搜索，并在最大可行速度附近采用二分法进行精细搜索，最终确定满足全部制程要求的最大可行速度。

对于最终结果，应进一步验证
\[
v_{\max}-\varepsilon
\]
仍属于可行域，而
\[
v_{\max}+\varepsilon
\]
不满足至少一项制程约束，从而证明所得到的速度确为可行域右端点。

\subsection{问题三结论}

问题三进一步将四组温区设定温度和传送带速度作为决策变量，以217℃以上至峰值之间的超温面积为目标函数：
\[
A
=
\int_{t_{217}^{\uparrow}}^{t_{pk}}
\left[T(t)-217\right]{\rm d}t.
\]

在峰值温度、升降温斜率、150--190℃升温时间和217℃以上持续时间等全部工艺约束下，构建非线性约束优化模型。

通过全局搜索与局部精化相结合的方法，在给定决策变量范围内寻找超温面积较小且满足全部工艺约束的温区温度组合和传送带速度。

最终应报告各温区最优设定温度、最优速度、超温面积以及全部制程指标，并通过独立程序重新计算以验证结果的可行性。

\subsection{问题四结论}

问题四在问题三超温面积优化的基础上进一步考虑峰值两侧曲线的形状对称性。

将峰值前后217℃以上区域分别映射至归一化时间区间 $\theta\in[0,1]$，定义
\[
D
=
\int_0^1
\left|
T_{\rm pre}(\theta)
-
T_{\rm post}(\theta)
\right|
{\rm d}\theta
\]
作为曲线对称性指标。

在此基础上，对超温面积与对称性指标进行无量纲化处理，并构造综合目标函数
\[
F
=
(1-\lambda)
\frac{A}{A_{\rm ref}}
+
\lambda
\frac{D}{D_{\rm ref}}.
\]

通过改变权重参数 $\lambda$ 进行敏感性分析，可以进一步观察最优方案在“减小超温面积”和“提高曲线对称性”之间的权衡关系。

\subsection{总体结论}

本文最终形成如下统一求解框架：
\[
\boxed{
\text{环境温度场}
\rightarrow
\text{温度响应}
\rightarrow
\text{制程指标}
\rightarrow
\text{可行域}
\rightarrow
\text{参数优化}
}
\]

该模型既能够描述给定温区条件下的炉温变化规律，也能够进一步用于传送带速度和温区设定温度的优化，为回焊炉工艺参数设计提供了一种具有物理解释性的数学方法。


% =========================================================
% 十二、参考文献
% =========================================================

\begin{thebibliography}{99}

\bibitem{ref1}
全国大学生数学建模竞赛组委会.
2020年高教社杯全国大学生数学建模竞赛A题：炉温曲线[EB/OL].
2020.

\bibitem{ref2}
杨世铭, 陶文铨.
传热学[M].
北京: 高等教育出版社.

\bibitem{ref3}
司守奎, 孙兆亮.
数学建模算法与应用[M].
北京: 国防工业出版社.

\bibitem{ref4}
Storn R, Price K.
Differential Evolution---A Simple and Efficient Heuristic for Global Optimization over Continuous Spaces[J].
Journal of Global Optimization, 1997, 11(4): 341--359.

\bibitem{ref5}
Nocedal J, Wright S J.
Numerical Optimization[M].
New York: Springer, 2006.

\end{thebibliography}


% =========================================================
% 附录
% =========================================================

\begin{appendices}

\section{问题一求解核心程序}

\subsection{环境温度函数}

\begin{verbatim}
def ambient_temperature(x, zone_temp):

    T1_5, T6, T7, T8_9, T10_11 = zone_temp

    Tset = [
        T1_5, T1_5, T1_5, T1_5, T1_5,
        T6, T7, T8_9, T8_9, T10_11, T10_11
    ]

    L = 30.5
    gap = 5.0
    front = 25.0
    rear = 25.0

    # 炉前区域
    if 0 <= x < front:
        return 25.0 + x / front * (Tset[0] - 25.0)

    # 各温区
    for i in range(11):

        xs = front + i * (L + gap)
        xe = xs + L

        if xs <= x <= xe:
            return Tset[i]

        # 温区间隙
        if i < 10:

            next_xs = front + (i + 1) * (L + gap)

            if xe < x < next_xs:

                ratio = (x - xe) / (next_xs - xe)

                return (
                    Tset[i]
                    + ratio * (Tset[i + 1] - Tset[i])
                )

    # 炉后冷却区域
    x11_end = front + 10 * (L + gap) + L

    if x11_end < x <= x11_end + rear:

        ratio = (x - x11_end) / rear

        return (
            Tset[10]
            + ratio * (25.0 - Tset[10])
        )

    return 25.0
\end{verbatim}


\subsection{温度微分方程}

\begin{verbatim}
def temperature_rhs(t, T, alpha, speed, zone_temp):

    x = speed * t

    Ta = ambient_temperature(
        x,
        zone_temp
    )

    return alpha * (Ta - T)
\end{verbatim}


\subsection{四阶Runge--Kutta数值积分}

\begin{verbatim}
def rk4_step(
    t,
    T,
    dt,
    alpha,
    speed,
    zone_temp
):

    k1 = temperature_rhs(
        t,
        T,
        alpha,
        speed,
        zone_temp
    )

    k2 = temperature_rhs(
        t + dt / 2,
        T + dt * k1 / 2,
        alpha,
        speed,
        zone_temp
    )

    k3 = temperature_rhs(
        t + dt / 2,
        T + dt * k2 / 2,
        alpha,
        speed,
        zone_temp
    )

    k4 = temperature_rhs(
        t + dt,
        T + dt * k3,
        alpha,
        speed,
        zone_temp
    )

    return (
        T
        + dt / 6
        * (
            k1
            + 2 * k2
            + 2 * k3
            + k4
        )
    )
\end{verbatim}


\subsection{综合换热参数辨识}

\begin{verbatim}
def calibration_error(
    alpha,
    t_obs,
    T_obs,
    zone_temp,
    speed
):

    t_model, T_model = simulate_temperature(
        alpha,
        speed,
        zone_temp
    )

    T_fit = interpolate(
        t_model,
        T_model,
        t_obs
    )

    return np.sum(
        (T_fit - T_obs) ** 2
    )


def calibrate_alpha(
    t_obs,
    T_obs,
    zone_temp,
    speed
):

    result = minimize_scalar(
        calibration_error,
        bounds=(1e-5, 1.0),
        method='bounded',
        args=(
            t_obs,
            T_obs,
            zone_temp,
            speed
        )
    )

    return result.x
\end{verbatim}


% =========================================================
% 附录B
% =========================================================

\section{问题二求解核心程序}

\subsection{制程约束判断}

\begin{verbatim}
def feasible_problem2(speed, alpha):

    zone_temp = [
        182.0,
        203.0,
        237.0,
        254.0,
        25.0
    ]

    t, T = simulate_temperature(
        alpha,
        speed,
        zone_temp
    )

    metrics = process_metrics(
        t,
        T
    )

    return (
        240.0 <= metrics["T_peak"] <= 250.0
        and
        0.0 <= metrics["S_r"] <= 3.0
        and
        -3.0 <= metrics["S_d"] <= 0.0
        and
        60.0 <= metrics["delta_150_190"] <= 120.0
        and
        40.0 <= metrics["delta_217"] <= 90.0
    )
\end{verbatim}


\subsection{最大允许速度搜索}

\begin{verbatim}
def solve_problem2(alpha):

    feasible = []

    # 第一步：粗搜索
    for speed in np.arange(
        65.0,
        100.0 + 1.0,
        1.0
    ):

        if feasible_problem2(
            speed,
            alpha
        ):
            feasible.append(speed)

    if len(feasible) == 0:
        return None

    # 第二步：确定边界区间
    v_left = max(feasible)
    v_right = min(
        v_left + 1.0,
        100.0
    )

    # 第三步：二分
    eps = 1e-3

    while v_right - v_left > eps:

        v_mid = (
            v_left
            + v_right
        ) / 2.0

        if feasible_problem2(
            v_mid,
            alpha
        ):
            v_left = v_mid
        else:
            v_right = v_mid

    return v_left
\end{verbatim}


\subsection{最大速度边界验证}

\begin{verbatim}
def verify_speed_boundary(
    alpha,
    v_star,
    eps=0.01
):

    v_minus = max(
        65.0,
        v_star - eps
    )

    v_plus = min(
        100.0,
        v_star + eps
    )

    return {
        "v_minus":
            feasible_problem2(
                v_minus,
                alpha
            ),

        "v_plus":
            feasible_problem2(
                v_plus,
                alpha
            )
    }
\end{verbatim}


% =========================================================
% 附录C
% =========================================================

\section{问题三求解核心程序}

\subsection{超温面积计算}

\begin{verbatim}
def excess_area(t, T):

    peak_index = np.argmax(T)

    t_up = t[:peak_index + 1]
    T_up = T[:peak_index + 1]

    above = np.where(
        T_up >= 217.0
    )[0]

    if len(above) == 0:
        return 0.0

    first = above[0]

    t217 = interpolate_crossing(
        t_up,
        T_up,
        217.0,
        first
    )

    mask = (
        (t_up >= t217)
        &
        (t_up <= t_up[peak_index])
    )

    return np.trapz(
        T_up[mask] - 217.0,
        t_up[mask]
    )
\end{verbatim}


\subsection{问题三目标函数}

\begin{verbatim}
def objective_problem3(
    x,
    alpha
):

    T1_5 = x[0]
    T6 = x[1]
    T7 = x[2]
    T8_9 = x[3]
    speed = x[4]

    zone_temp = [
        T1_5,
        T6,
        T7,
        T8_9,
        25.0
    ]

    t, T = simulate_temperature(
        alpha,
        speed,
        zone_temp
    )

    metrics = process_metrics(
        t,
        T
    )

    violation = constraint_violation(
        metrics
    )

    area = excess_area(
        t,
        T
    )

    if violation <= 1e-8:
        return area

    return (
        area
        + 1e6 * violation
    )
\end{verbatim}


\subsection{问题三优化}

\begin{verbatim}
def solve_problem3(alpha):

    bounds = [
        (165.0, 185.0),
        (185.0, 205.0),
        (225.0, 245.0),
        (245.0, 265.0),
        (65.0, 100.0)
    ]

    result = differential_evolution(
        lambda x:
            objective_problem3(
                x,
                alpha
            ),
        bounds=bounds,
        strategy='best1bin',
        maxiter=500,
        popsize=20,
        tol=1e-7,
        polish=True,
        seed=2026
    )

    return result
\end{verbatim}


% =========================================================
% 附录D
% =========================================================

\section{问题四求解核心程序}

\subsection{峰值两侧曲线归一化}

\begin{verbatim}
def symmetry_index(
    t,
    T,
    n_grid=500
):

    peak_index = np.argmax(T)

    t_peak = t[peak_index]

    above = np.where(
        T > 217.0
    )[0]

    if len(above) < 2:
        return np.inf

    first = above[0]
    last = above[-1]

    t_in = interpolate_crossing(
        t,
        T,
        217.0,
        first
    )

    t_out = interpolate_crossing(
        t,
        T,
        217.0,
        last
    )

    theta = np.linspace(
        0.0,
        1.0,
        n_grid
    )

    t_pre = (
        t_peak
        - theta
        * (t_peak - t_in)
    )

    t_post = (
        t_peak
        + theta
        * (t_out - t_peak)
    )

    T_pre = interpolate(
        t,
        T,
        t_pre
    )

    T_post = interpolate(
        t,
        T,
        t_post
    )

    return np.trapz(
        np.abs(
            T_pre - T_post
        ),
        theta
    )
\end{verbatim}


\subsection{问题四综合目标函数}

\begin{verbatim}
def objective_problem4(
    x,
    alpha,
    A_ref,
    D_ref,
    lam=0.4
):

    zone_temp = [
        x[0],
        x[1],
        x[2],
        x[3],
        25.0
    ]

    speed = x[4]

    t, T = simulate_temperature(
        alpha,
        speed,
        zone_temp
    )

    metrics = process_metrics(
        t,
        T
    )

    violation = constraint_violation(
        metrics
    )

    A = excess_area(
        t,
        T
    )

    D = symmetry_index(
        t,
        T
    )

    if violation > 1e-8:

        return (
            1e6 * violation
            +
            A / max(A_ref, 1e-8)
            +
            D / max(D_ref, 1e-8)
        )

    return (
        (1.0 - lam)
        * A / A_ref
        +
        lam
        * D / D_ref
    )
\end{verbatim}


\subsection{问题四优化}

\begin{verbatim}
def solve_problem4(
    alpha,
    A_ref,
    D_ref
):

    bounds = [
        (165.0, 185.0),
        (185.0, 205.0),
        (225.0, 245.0),
        (245.0, 265.0),
        (65.0, 100.0)
    ]

    result = differential_evolution(
        lambda x:
            objective_problem4(
                x,
                alpha,
                A_ref,
                D_ref,
                0.4
            ),
        bounds=bounds,
        strategy='best1bin',
        maxiter=500,
        popsize=20,
        tol=1e-7,
        polish=True,
        seed=2026
    )

    return result
\end{verbatim}


% =========================================================
% 附录E
% =========================================================

\section{数值计算与结果复核}

\subsection{步长敏感性分析}

\begin{verbatim}
def sensitivity_analysis(
    alpha,
    speed,
    zone_temp
):

    dt_list = [
        1.0,
        0.5,
        0.25,
        0.125
    ]

    results = {}

    for dt in dt_list:

        t, T = simulate_temperature(
            alpha,
            speed,
            zone_temp,
            dt
        )

        results[dt] = (
            process_metrics(
                t,
                T
            )
        )

    return results
\end{verbatim}


\subsection{参数敏感性分析}

\begin{verbatim}
def alpha_sensitivity(
    alpha_star,
    speed,
    zone_temp
):

    perturbations = [
        -0.10,
        -0.05,
         0.00,
         0.05,
         0.10
    ]

    results = {}

    for eta in perturbations:

        alpha = (
            1.0 + eta
        ) * alpha_star

        t, T = simulate_temperature(
            alpha,
            speed,
            zone_temp
        )

        results[eta] = (
            process_metrics(
                t,
                T
            )
        )

    return results
\end{verbatim}


\subsection{最终制程约束检查}

\begin{verbatim}
def final_check(metrics):

    c1 = (
        240.0
        <= metrics["T_peak"]
        <= 250.0
    )

    c2 = (
        0.0
        <= metrics["S_r"]
        <= 3.0
    )

    c3 = (
        -3.0
        <= metrics["S_d"]
        <= 0.0
    )

    c4 = (
        60.0
        <= metrics["delta_150_190"]
        <= 120.0
    )

    c5 = (
        40.0
        <= metrics["delta_217"]
        <= 90.0
    )

    return (
        c1
        and c2
        and c3
        and c4
        and c5
    )
\end{verbatim}


\subsection{结果记录}

最终计算结果统一记录为：

\begin{table}[H]
\centering
\caption{最终优化方案结果记录}
\begin{tabular}{ccc}
\toprule
变量或指标 & 计算结果 & 单位\\
\midrule
综合换热参数 $\alpha$ & 待程序填入 & s$^{-1}$\\
最大允许速度 $v_{\max}$ & 待程序填入 & cm/min\\
问题一位置温度 & 待程序填入 & ${}^\circ$C\\
问题三超温面积 $A$ & 待程序填入 & ${}^\circ$C$\cdot$s\\
问题四对称性指标 $D$ & 待程序填入 & ${}^\circ$C\\
\bottomrule
\end{tabular}
\end{table}

\end{appendices}

\end{document}